\begin{recipe}
[% 
    preparationtime = {\unit[25]{min}},
    bakingtime = {\unit[5--8]{min}},
    bakingtemperature = {\protect\bakingtemperature{
        topbottomheat=\unit[200]{\textcelcius}
    }},
    portion = {\portion{2}},
    %source = {},
    calory = {240kcal | 3g Protein | 33g KH | 12g Fett},
    price = {1,60€},
    created = {11.06.2026},
    %rating = {1},
]
{Reispapier-Pillows mit Guacamole}

\graph{% pictures
    big=pic/vorspeise/06_reispapier-pillows-guacamole
}

\introduction{%
Knusprig aufgepuffte Reispapier-Kissen, gefüllt mit cremiger Guacamole. Sieht deutlich komplizierter aus als es ist und ist damit genau die richtige Art von Vorspeise.
}

\ingredients{
    \textbf{Pillows} & \\
    6 Blätter & Reispapier\index{Getreide!Reispapier}\\
    & Wasser zum Anfeuchten\\
    \textbf{Guacamole} & \\
    1 & Reife Avocado\index{Obst!Avocado}\\
    \nicefrac{1}{2} & Tomate\\
    \nicefrac{1}{2} & Limette\\
    \nicefrac{1}{4} & Rote Zwiebel\\
    Optional & Koriander oder Petersilie\\
    & Salz, Pfeffer\\
    & Chiliflocken
}

\preparation{%
    \step%
    Backofen auf 200 °C Ober-/Unterhitze vorheizen. Ein Backblech mit Backpapier belegen oder ein Gitter auf ein Backblech setzen.
    
    \step%
    Je zwei Reispapierblätter kurz mit Wasser anfeuchten und direkt aufeinanderlegen. Sie sollen weich genug sein, um zusammenzukleben, aber nicht komplett matschig werden.
    
    \step%
    Die doppelten Reispapierblätter in Quadrate schneiden und mit etwas Abstand auf das Blech oder Gitter legen.
    
    \step%
    Im Ofen 5--8 Minuten backen, bis die Stücke aufpuffen und kleine Kissen entstehen. Dabei gut beobachten, weil Reispapier ziemlich schnell von perfekt zu verbrannt wechseln kann.
    
    \step%
    Die Pillows vollständig auskühlen lassen. Danach mit einem heißen Metallspieß, einer Messerspitze oder vorsichtig mit einer Schere ein kleines Loch hineinmachen.
    
    \step%
    Für die Guacamole Avocado zerdrücken. Tomate entkernen und fein würfeln. Zwiebel sehr fein schneiden.
    
    \step%
    Avocado mit Tomate, Zwiebel, Limettensaft, Salz, Pfeffer und Chiliflocken vermengen. Optional etwas Koriander oder Petersilie unterheben.
    
    \step%
    Guacamole in einen Spritzbeutel oder Gefrierbeutel mit abgeschnittener Ecke geben und vorsichtig in die Reispapier-Pillows füllen.
    
    \step%
    Direkt servieren, solange die Pillows noch knusprig sind.
}

\hint{
    Erst kurz vor dem Servieren füllen, sonst zieht die Guacamole Feuchtigkeit ins Reispapier und die Pillows werden wieder weich. Also leider kein Meal Prep, sondern eher kleines Show-off-Ding.
}

\suggestion[Variation]{%
    Für eine etwas frischere Füllung kann man die Guacamole mit fein gewürfelter Gurke strecken. Dann aber wirklich gut abtropfen lassen, sonst killt die Feuchtigkeit den Crunch.
}

\end{recipe}